\thesisfrontchapter{Abstract}

% PolyU requires 200--500 words. Replace this example before submission.
This thesis investigates a clearly defined research problem whose practical importance and unresolved scientific questions motivate the work. It begins by reviewing the relevant theoretical foundations and empirical evidence, identifying limitations in existing approaches and establishing the research questions addressed in the subsequent chapters. A reproducible methodology is then developed to examine these questions using appropriate data, experimental controls, evaluation measures, and statistical analyses. The proposed approach is compared with established baselines under consistent conditions, while additional experiments examine robustness, sensitivity to design choices, and the influence of the principal assumptions. The results demonstrate how the selected methodology addresses the stated research objectives and clarify the circumstances in which its advantages and limitations become most apparent. The discussion relates the findings to prior research, considers alternative explanations, and distinguishes conclusions supported directly by evidence from interpretations that require further validation. Particular attention is given to reproducibility, uncertainty, generalisability, ethical considerations, and the practical implications of the work. Taken together, the studies provide a coherent contribution to the field and establish a foundation for subsequent investigation. Future research should extend the evaluation to additional settings, test the approach using independent evidence, and examine how the identified limitations can be mitigated. Replace this anonymous example with a self-contained summary of the completed thesis, including its actual problem, methods, principal quantitative findings, conclusions, and significance.
